\documentclass[10pt,letterpaper]{article}
\usepackage[margin=0.7in,headheight=14pt]{geometry}
\usepackage[T1]{fontenc}
\usepackage{lmodern,amsmath,booktabs,xcolor,fancyhdr,microtype}
\usepackage[colorlinks=true,urlcolor=teal]{hyperref}
\definecolor{navy}{HTML}{192B4D}
\definecolor{teal}{HTML}{007D80}
\pagestyle{fancy}\fancyhf{}
\fancyhead[L]{\small MF 731 | Homework 5}
\fancyhead[R]{\small 09: Stress Testing}
\fancyfoot[C]{\small\thepage}
\setlength{\parindent}{0pt}
\setlength{\parskip}{6pt}
\hypersetup{pdftitle={MF 731 Homework 5 - Stress Testing}}
\begin{document}
{\Large\bfseries\textcolor{navy}{Stress Testing a Microsoft--Apple Portfolio}}

\textbf{Data and EWMA.}
I use Yahoo Finance daily closing prices, adjusted for stock splits but not
dividends. The 1,259 log returns are dated September 1, 2011 through
September 1, 2016; the August 31, 2011 close supplies the first return.
The given market caps imply weights $(0.438303,0.561697)$, so the
\$1,000,000 portfolio has fixed initial dollar positions
$\theta_t=(438{,}303.32,561{,}696.68)^\top$, ordered as (MSFT, AAPL).
Starting from zero, I update the mean and covariance with decay $0.97$,
using the previous mean in each covariance residual. The final estimates are
\[
\mu=\begin{pmatrix}0.00125025\\0.00106069\end{pmatrix},\qquad
\Sigma=\begin{pmatrix}
0.000129897&0.000033007\\
0.000033007&0.000151852
\end{pmatrix},\qquad \rho=0.235013.
\]

\textbf{Initial VaR and capital.}
For the normally distributed one-day linearized loss,
\[
\mathrm{VaR}_{0.95,1}
=-\theta_t^\top\mu+1.644854\sqrt{\theta_t^\top\Sigma\theta_t}
=\$14{,}383.87.
\]
Applying the requested square-root-of-time rule gives
\[
\mathrm{VaR}_{0.95,10}^{\mathrm{scale}}
=\sqrt{10}\,\mathrm{VaR}_{0.95,1}
=\boxed{\$45{,}485.80},\qquad
C=3\,\mathrm{VaR}_{0.95,10}^{\mathrm{scale}}
=\boxed{\$136{,}457.39}.
\]

\textbf{Shock and simulation.}
Apple's five-sigma log-return shock is $x^{(2)}=-0.0605534$.
Conditional on this value, Microsoft's mean and standard deviation are
$\widetilde\mu^{(1)}=-0.0121422$ and
$\widetilde\sigma^{(1)}=0.0110780$. Its conditional five-sigma shock is
\[
x^{(1)}=\widetilde\mu^{(1)}-5\widetilde\sigma^{(1)}=-0.0675324.
\]
Both first-day returns are fixed. After updating EWMA with this shock,
I simulate days 2--10 from the current bivariate normal distribution,
updating the mean and covariance separately along each path.
For each of 50,000 paths (seed 731), I compute
\[
\ell_m=-\theta_t^\top\sum_{k=1}^{10}X_{t+k\Delta}^{m}.
\]
The empirical 95\% VaR is the 47,500th sorted loss, without interpolation.

\begin{center}
\begin{tabular}{llr}\toprule
 & Requested result & Estimate\\\midrule
(a) & Average stressed ten-day loss & \$70,838.04\\
(b) & Stressed ten-day 95\% VaR & \$151,692.60\\
(c) & Loss above initial ten-day VaR & 69.91\% \; (34,955/50,000)\\
(d) & Loss above initial regulatory capital & 8.89\% \; (4,445/50,000)\\\bottomrule
\end{tabular}
\end{center}

The stressed VaR is about 3.34 times the initial ten-day VaR and exceeds
the initial capital by \$15,235.21. Under this stress scenario, the capital
is exceeded more often than the 5\% rate associated with a 95\% VaR.
As a simulation check, the exact conditional mean loss is \$70,802.45;
the simulated mean is only 0.16 Monte Carlo standard errors away.

\vfill
\textbf{Python code, data and running instructions:}\\
\url{https://github.com/jesuswalking/MF-731-HW1/tree/main/HW5}
\end{document}
